\chapter{Vaginal Discharge}

\section*{Definition}
Vaginal discharge is fluid that comes from your vagina. Some discharge is normal and healthy, but changes in color, odor, or texture can signal an infection or other problem needing evaluation.

\section*{What Happens in Your Body}
Normal discharge is clear or white with little smell and changes with your menstrual cycle. Abnormal discharge may be yellow-green, gray, frothy, thick and clumpy, or bloody, often with itching, burning, or pain. Causes include bacterial vaginosis, yeast infections, sexually transmitted infections, and other conditions.

\section*{Causes}
\begin{itemize}
  \item Bacterial vaginosis
  \item Yeast infections Candida
  \item Trichomoniasis sexually transmitted
  \item Cervicitis STDs
  \item Atrophic vaginitis postmenopausal
\end{itemize}

\section*{When to See a Doctor}
See your doctor if:
\begin{itemize}
  \item Symptoms are severe or getting worse over time
  \item Discharge is new, has a bad odor, changed color, or is accompanied by itching or pain
  \item You have other concerning symptoms such as fever, weight loss, or bleeding
\end{itemize}

\section*{Related Symptoms}
\begin{itemize}
  \item Vaginal itching burning
  \item Pelvic pain abnormal bleeding
  \item pain during intercourse painful urination
  \item Lower abdominal pain
  \item Urinary frequency
\end{itemize}
\textbf{Important:} If this symptom persists, your doctor will check for serious underlying causes.

\section*{Self-Care / Home Management}
\begin{itemize}
  \item Track symptoms including cycle dates, flow characteristics, and associated pain using a diary or app.
  \item Apply a heating pad to the lower abdomen for menstrual cramp relief.
  \item Practice good genital hygiene and avoid douching or scented products which can disrupt normal flora.
  \item Over-the-counter pain relievers (ibuprofen, naproxen) are effective for menstrual pain and can reduce heavy bleeding.
  \item Seek evaluation for abnormal bleeding, pelvic pain, unusual discharge, or breast changes.
\end{itemize}

\section*{How Your Doctor Will Evaluate You}
\begin{itemize}
  \item History includes menstrual history, sexual activity, contraceptive use, pregnancy status, and associated symptoms.
  \item Pelvic examination (speculum and bimanual) assesses the cervix, uterus, and adnexa for tenderness, masses, or discharge.
  \item Pregnancy testing is essential for any reproductive-age woman with abnormal bleeding or pelvic symptoms.
  \item Ultrasound (transabdominal or transvaginal) evaluates the endometrium, myometrium, ovaries, and adnexal structures.
\end{itemize}

\section*{Treatment Options}
\begin{itemize}
  \item Hormonal contraceptives (oral, injectable, intrauterine) are first-line treatment for many menstrual disorders.
  \item NSAIDs reduce chemical messengers that cause inflammation and pain-mediated pain and bleeding in dysmenorrhea and menorrhagia.
  \item Specific diagnoses require targeted therapy: antibiotics for pelvic infection, endometrial ablation for refractory bleeding, surgery for fibroids or endometriosis.
  \item Referral to gynecology is indicated for complex cases, infertility, or suspected malignancy.
\end{itemize}

\section*{References}
\begin{thebibliography}{99}
\bibitem{vaginal_discharge:ref1} Workowski KA, Bachmann LH, Chan PA, et al. ``Sexually Transmitted Infections Treatment Guidelines, 2021.'' \textit{MMWR Recommendations and Reports}, 2021;70(4):1--187.
\bibitem{vaginal_discharge:ref2} Sobel JD. ``Vulvovaginal candidosis.'' \textit{The Lancet}, 2007;369(9577):1961--1971.
\bibitem{vaginal_discharge:ref3} Sobel JD. ``Bacterial vaginosis.'' \textit{Annual Review of Medicine}, 2000;51:349--356.
\bibitem{vaginal_discharge:ref4} Sonnex C. ``Trichomonas vaginalis infection in women: Diagnosis and treatment.'' \textit{British Medical Journal}, 2008;336(7657):1221--1224.
\bibitem{vaginal_discharge:ref5} Schwebke JR, Gaydos CA, Nyirjesy P, et al. ``Diagnostic performance of a molecular test versus clinician assessment of vaginal discharge.'' \textit{Journal of Clinical Microbiology}, 2018;56(6):e00222--18.
\bibitem{vaginal_discharge:ref6} Nyirjesy P. ``Management of persistent vaginitis.'' \textit{Obstetrics \& Gynecology}, 2014;124(6):1135--1146.
\end{thebibliography}
